\section{The dimension-drop criterion without stable rank one}
\label{sec:rw}

This section proves the implication of R\o rdam and Winter
\cite[Proposition~5.1, (ii)$\Rightarrow$(iv)]{RW} used in the finite
construction. We retain their order-zero and flip-path method, while
making the support projections and endpoint domains explicit. The
argument applies to an arbitrary unital $C^*$-algebra.

For $b\ge0$, write $\operatorname{Her}(b)=\overline{bBb}$.
Support projections and polar decompositions will be taken in $B^{**}$.
They need not belong to $B$; we will verify membership in $B$ whenever
they are used to construct an element of the original algebra.

\begin{lemma}\label{lem:rw-exact-cut}
If $a,b\in B_+$ satisfy $a\precsim b$, then for every $\delta>0$
there is $z\in B$ such that
\[
 z^*z=(a-\delta)_+,
 \qquad zz^*\in\operatorname{Her}(b).
\]
\end{lemma}
\begin{proof}
Since $a,b$ lie in $B$, compression of stabilized comparison
witnesses to their $(1,1)$ entries gives witnesses in $B$ itself.
Choose $x\in B$ with $\|a-x^*bx\|<\eta<\delta$, and put
$e=(a-\delta)_+$ and $y=b^{1/2}xe^{1/2}$. Then
\[
 T=y^*y\ge e^{1/2}(a-\eta)e^{1/2}
       \ge(\delta-\eta)e,
 \qquad T\le\|x^*bx\|e.
\]
Thus $T$ and $e$ have the same support. If $y=uT^{1/2}$ is the polar
decomposition, set $z=ue^{1/2}$ in $B^{**}$. To see that $z\in B$,
approximate it by
\[
 z_\tau=y(T+\tau)^{-1/2}e^{1/2}\in b^{1/2}B,
 \qquad \tau>0.
\]
Writing $f_\tau(t)=1-\sqrt{t/(t+\tau)}$, the domination of $e$ by
$T$ gives
\[
 \|z-z_\tau\|^2
 \le\|f_\tau(T)e f_\tau(T)\|
 \le\frac{\sup_{t\ge0}t f_\tau(t)^2}{\delta-\eta}
 \le\frac{\tau}{\delta-\eta}.
\]
No commutation of $e$ with $T$ is used here. Hence $z\in B$,
$z^*z=e$, and $zz^*$ is a norm limit of elements of
$b^{1/2}Bb^{1/2}$, as required.
\end{proof}

We shall use a concrete description of order-zero maps. If
$x_1,\ldots,x_q$ are contractions with
$x_i x_j^*=\delta_{ij}h$ for a positive contraction $h$, then
\[
 \alpha(e_{ij})=x_i^*x_j
\]
defines a completely positive contractive order-zero map from $M_q$.
Indeed, complete positivity follows from this Gram formula. The
positive elements $x_i^*x_i$ are orthogonal, so their sum $H$ is a
contraction; multiplication gives
$\alpha(e_{ij})\alpha(e_{kl})=\delta_{jk}H\alpha(e_{il})$.
This identity proves preservation of orthogonality. It also gives a
supporting homomorphism $\bar\alpha:M_q\to B^{**}$ with
\[
 \alpha(x)=H\bar\alpha(x)=\bar\alpha(x)H,
 \qquad \bar\alpha(1_q)=\operatorname{supp}(H).
\]
For completeness, its matrix units are the strong limits of
$\alpha(e_{ij})(H+\tau)^{-1}$ as $\tau\downarrow0$. The displayed
multiplication identity and
$\alpha(e_{ij})^*\alpha(e_{ij})\le H^2$ verify their boundedness,
matrix-unit relations, and stated support.

\begin{lemma}\label{lem:rw-commuting-maps}
Let $\alpha:M_3\to B$ and $\beta:M_2\to B$ be completely positive
contractive maps with commuting ranges, satisfying
$\psi(x)\psi(y)=\psi(1)\psi(xy)$ for $\psi=\alpha,\beta$.
If
$\alpha(1_3)+\beta(1_2)=1_B$, there is a unital homomorphism
$I(2,3)\to B$.
\end{lemma}
\begin{proof}
This is the order-zero presentation used in
\cite[Proposition~2.5]{RW}; we include the verification.
The canonical maps into $I(2,3)$ are
\[
 \alpha_0(y)(t)=t(1_2\otimes y),\qquad
 \beta_0(x)(t)=(1-t)(x\otimes1_3).
\]
Their generated algebra contains the scalar function $t$, and products
of their entries generate the ideal $C_0((0,1),M_2\otimes M_3)$.
Their endpoint values also generate $M_2\oplus M_3$. Consequently
they generate all of $I(2,3)$.

Put $E=C^*(\alpha(M_3),\beta(M_2),1_B)$. The element
$\alpha(1_3)$ is central in $E$. In any irreducible representation
of $E$ it is a scalar $r\in[0,1]$. If $0<r<1$, the rescaled maps
$r^{-1}\alpha$ and $(1-r)^{-1}\beta$ in this representation are
commuting unital matrix homomorphisms, and hence define a representation
of $M_2\otimes M_3$. At $r=0$ only the $M_2$ map remains, and at
$r=1$ only the $M_3$ map remains. These assertions follow either from
the supporting homomorphisms or directly from the order-zero
multiplication identity.

It follows that substitution of $\alpha,\beta$ for
$\alpha_0,\beta_0$ in a noncommutative $*$-polynomial is bounded in
norm by the norm of the polynomial in $I(2,3)$: test in each
irreducible representation of $E$ and use the corresponding evaluation
at $r$. Substitution is therefore well defined and extends from the
dense generated algebra to the required unital homomorphism.
\end{proof}

\begin{proof}[Proof of Proposition~\ref{prop:rw}]
\emph{Exact relations from the comparison.}
Write $a_1=w^*w$ and $a_2=ww^*$ for the equivalence witness.
Decreasing $\varepsilon$ if necessary, assume $0<\varepsilon<1$.
Let $f(t)=\min\{t/\varepsilon,1\}$ and put
$s_1=f(a_1)^{1/2}$. If $w^*=u_0a_2^{1/2}$ is its polar
decomposition, put $s_2=u_0f(a_2)^{1/2}$. This belongs to $B$:
it equals $w^*q(a_2)$, where the continuous function
$q(t)=\sqrt{f(t)/t}$ for $t>0$ is assigned the value
$\varepsilon^{-1/2}$ at zero. Thus
\[
 s_i s_j^*=\delta_{ij}s_1^2,\qquad
 d:=s_1^*s_1+s_2^*s_2=f(a_1)+f(a_2).
\]
Here the off-diagonal products vanish because the initial positive
elements are orthogonal. Writing $c=1-a_1-a_2$, functional calculus
in the commuting algebra generated by $a_1,a_2$ gives
\[
 a:=1-d=\varepsilon^{-1}(c-(1-\varepsilon))_+.
\]
Apply Lemma~\ref{lem:rw-exact-cut} to $c\precsim(a_1-\varepsilon)_+$,
with $\delta=1-\varepsilon$, and divide its element by
$\sqrt\varepsilon$. We obtain a contraction $v\in B$ satisfying
\begin{equation}\label{eq:rw-relations}
 v^*v=a=1-d,\qquad vv^*\in\operatorname{Her}((a_1-\varepsilon)_+),
 \qquad s_1^2vv^*=vv^*.
\end{equation}
In particular $s_1v=v$, $s_2v=0$, and $av=0$.
The initial and final support relations also give
$ds_i=s_1^2s_i=s_id$; hence $a$ commutes with the $s_i$.
The Gram construction gives $\rho:M_2\to B$ with
$\rho(e_{ij})=s_i^*s_j$ and $\rho(1_2)=d$.
Write $\sigma$ for its supporting homomorphism and
$p=\operatorname{supp}(d)$.

\emph{Two orthogonal copies of the defect algebra.}
Let $v=ua^{1/2}$ be the polar decomposition and
$e=u^*u=\operatorname{supp}(a)$. Set
\[
 r_i=s_i^*u=\sigma(e_{i1})u\quad(i=1,2),
 \qquad R=r_1r_1^*+r_2r_2^*.
\]
The second formula for $r_i$ follows from $du=u$ and $s_1u=u$.
Direct multiplication gives
\begin{equation}\label{eq:rw-corners}
 r_i^*r_j=\delta_{ij}e,\qquad
 dr_i=r_i,\qquad R\le p,\qquad R\perp e,\qquad
 \sigma(e_{jk})r_i=\delta_{ki}r_j.
\end{equation}
For $R\perp e$, note that $a$ commutes with each $s_i$ (because
$d$ does) and $au=0$. The map
\[
 \varphi:\operatorname{Her}(a)\longrightarrow B,
 \qquad \varphi(b)=\sum_{i=1}^2 r_i b r_i^*
\]
is a homomorphism by \eqref{eq:rw-corners}. Its values belong to
$B$: approximate $b\in\operatorname{Her}(a)$ by
$a^{1/2}b_m a^{1/2}$ and use $ua^{1/2}=v$ to see that
$ubu^*\in B$. The same approximation, with one outer factor omitted,
shows that $ub\in B$.
The same relations show
\begin{equation}\label{eq:rw-phi}
 \varphi(b)r_i=r_i b,\qquad
 [\varphi(b),\sigma(M_2)]=0,\qquad
 \varphi(\operatorname{Her}(a))\perp\operatorname{Her}(a).
\end{equation}

Here $\varphi(b)=\sum_i r_i b r_i^*$ places the same defect element in
the two ranges of the $r_i$, while $\sigma(e_{jk})$ acts on their index
$i$. To obtain the other coordinate, we use sections on $[0,1]$ with values
in $M_2$, evaluated at $d$, that are scalar at zero and vanish at one.
The scalar value accounts for the corner outside $p=\operatorname{supp}(d)$;
vanishing at one keeps the evaluated section in $\operatorname{Her}(1-d)$.
Define
\[
 D_1=\{F\in C([0,1],M_2):F(0)\in\mathbb C1_2,\ F(1)=0\},
\]
and let $c_F$ be the scalar with $F(0)=c_F1_2$. Since $d$ commutes
with $\sigma(M_2)$, continuous functional calculus defines
$\bar\lambda(F)=\sum_{i,j}F_{ij}(d)\sigma(e_{ij})$ in $pB^{**}p$.
The required homomorphism is
\begin{equation}\label{eq:rw-lambda}
 \lambda(F)=\bar\lambda(F)+c_F(1-p).
\end{equation}
The extra summand retains the scalar fiber at zero. In particular
$\lambda(f1_2)=f(d)$ for scalar $f$ vanishing at one.
The algebra $D_1$ is generated by $(1-t)1_2$ and
$t(1-t)x$, $x\in M_2$. Their images are respectively $a$ and
$a\rho(x)$, so $\lambda$ takes values in $\operatorname{Her}(a)$.
This also verifies membership in $B$, without assuming $p\in B$.

Let
\[
 D_2=\{F\in C([0,1],M_2\otimes M_2):
                   F(0)\in\mathbb C1_4,\ F(1)=0\}.
\]
It is a subalgebra of $D_1\otimes M_2$. By
\eqref{eq:rw-phi}, the commuting product of
$\varphi\circ\lambda$ and $\sigma$ defines a homomorphism on
$D_1\otimes M_2$; denote its restriction to $D_2$ by $\gamma$.
Its range lies in $B$: since $d\varphi(b)=\varphi(b)$,
$\varphi(b)\sigma(x)=\varphi(b)\rho(x)\in B$.
Its support projection is $R$, because its range is supported there
and contains $\varphi(a)$, which has support $R$.

The product defining $\gamma$ pairs this transported section action with
the action of $\sigma$ on the copy index. A weighted flip will interchange
them only on the overlap needed for the final commutation calculation.

\emph{The flip where the two pieces overlap.}
Let $F\in M_2\otimes M_2$ be the flip unitary. Choose a continuous
function $\theta:[0,1]\to[0,1]$ that is zero on $[0,1/3]$ and one
on $[2/3,1]$, and set
\[
 U(t)=\frac{1+F}{2}+e^{i\pi\theta(t)}\frac{1-F}{2}.
\]
This is an ordinary unitary path, equal to $1$ near zero and to
$F$ on $[2/3,1]$. It is a multiplier of $D_2$. The multiplier
extension of $\gamma$ sends it to a unitary $w_0$ in the corner
$RB^{**}R$; put $W=w_0+1-R$.

Let $g(t)=1$ for $0\le t\le2/3$ and $g(t)=3(1-t)$ otherwise.
Put
\[
 h=\lambda(g1_2)=g(d),\qquad k=\varphi(h),\qquad l=1-h-k.
\]
The elements $h,k$ are orthogonal positive contractions, and
$g(t)\ge1-t$ implies $0\le l\le d$. Also $kW=Wk$ and
\[
 k^\alpha W=\gamma(g^\alpha U)\in B\qquad(\alpha>0).
\]
The formula is legitimate because $g^\alpha U\in D_2$.
Define
\[
 x_i=h^{1/2}r_i^*W\quad(i=1,2),\qquad x_3=h^{1/2}.
\]
These elements belong to $B$: for $i\le2$ use the factorization
\[
 x_i=(r_i^*k^{1/4})(k^{1/4}W),\qquad
 r_i^*k^{1/4}=h^{1/4}r_i^*\in B.
\]
By \eqref{eq:rw-corners}, $R\perp e$, and unitarity of $W$,
\[
 x_i x_j^*=\delta_{ij}h\quad(1\le i,j\le3),\qquad
 \sum_{i=1}^3x_i^*x_i=h+k.
\]
Thus $\alpha(e_{ij})=x_i^*x_j$ defines an order-zero map
$\alpha:M_3\to B$ with $\alpha(1_3)=h+k$.

Since $h,k$ commute with $\sigma(M_2)$ and $0\le l\le d$, the
formula $\beta(y)=l\sigma(y)$ defines an order-zero map
$\beta:M_2\to B$. To justify its range directly, $l$ vanishes on
the spectral subspace of $d$ for $[0,2/3]$, since $h=1$ there and
$hk=0$. Choose a continuous scalar function $j$ that equals $1/t$
for $t\ge2/3$. Then
$\beta(y)=lj(d)\rho(y)\in B$.
Since $lp=l$, we have
$\alpha(1_3)+\beta(1_2)=1$.

It remains to check that their ranges commute. This is precisely
where the flip is needed. The scalar function $Q=g(1-g)$ vanishes
at both endpoints and is supported where $U=F$. Therefore
\[
 Q(t)(1\otimes y)U(t)^*
   =U(t)^*Q(t)(y\otimes1),\qquad y\in M_2.
\]
Both sides are elements of $D_2$. Applying $\gamma$ gives the
well-defined identity
\begin{equation}\label{eq:rw-weighted-flip}
 k(1-k)\sigma(y)W^*
       =W^*\varphi\bigl(\lambda(Q\otimes y)\bigr).
\end{equation}
In particular, no constant nonscalar section is evaluated by
$\lambda$. For $i=1,2$, put
$C_i=\alpha(e_{i3})=W^*r_i h$. Using
\eqref{eq:rw-phi} and \eqref{eq:rw-weighted-flip}, we obtain
\[
\begin{aligned}
 \beta(y)C_i
 &=k(1-k)\sigma(y)W^*r_i\\
 &=W^*r_i\lambda(Q\otimes y)\\
 &=W^*r_i h(1-h)\sigma(y)
  =C_i\beta(y).
\end{aligned}
\]
Here $\lambda(Q\otimes y)=Q(d)\sigma(y)$ because $Q(0)=0$;
the last equality uses $hk=0$.

The $C^*$-algebra generated by the $C_i$ contains $\alpha(M_3)$. Indeed,
$C_i^*C_i=h^2$ and $C_iC_i^*=\alpha(e_{ii})^2$, so their generated
algebra contains $h$, all diagonal entries, and
$A_0:=\alpha(1_3)$. Moreover
$C_iC_j^*=A_0\alpha(e_{ij})$, and multiplication by
$(A_0+\tau)^{-1}$ followed by $\tau\downarrow0$ recovers every
remaining matrix entry. The convergence follows from
$\alpha(e_{ij})=A_0\bar\alpha(e_{ij})$ and
$\|\tau A_0(A_0+\tau)^{-1}\|\le\tau$.
Thus $\alpha(M_3)$ commutes with $\beta(M_2)$.
Lemma~\ref{lem:rw-commuting-maps} completes the proof.

The construction also covers $v=0$: then $a=h=k=0$, $R=0$,
$W=1$, and $\beta$ is a unital $M_2$ homomorphism. No step requires
stable rank one or a nonzero defect.
\end{proof}
